\textbf{Phần 1. Trắc nghiệm lựa chọn}
\setcounterpageref{ex}{0}

\begin{ex} %[5D4V19]
Tính giá trị biểu thức $A = 14{,}25 + 8{,}6 + 5{,}75 + 1{,}4$ bằng cách thuận tiện nhất:
\choice
{$29$}
{\True $30$}
{$31$}
{$30{,}5$}
\loigiai{
Nhóm các số hạng có tổng là số tự nhiên:
\[
A = (14{,}25 + 5{,}75) + (8{,}6 + 1{,}4) = 20 + 10 = 30.
\]
\textbf{Đáp án của câu hỏi là B}
}
\end{ex}

\begin{ex} %[5D4V19]
Khi cộng một số tự nhiên với một số thập phân có một chữ số ở phần thập phân, bạn Nam quên viết dấu phẩy của số thập phân nên đã thực hiện cộng như hai số tự nhiên và được kết quả là $218$. Biết kết quả đúng là $50{,}6$. Số tự nhiên đó là
\choice
{$31$}
{\True $32$}
{$33$}
{$34$}
\loigiai{
Khi quên viết dấu phẩy của số thập phân có một chữ số ở phần thập phân, số đó đã tăng lên $10$ lần.\\
Hiệu giữa tổng sai và tổng đúng là:
\[
218 - 50{,}6 = 167{,}4.
\]
Số thập phân đó là:
\[
167{,}4 : (10 - 1) = 167{,}4 : 9 = 18{,}6.
\]
Số tự nhiên ban đầu là:
\[
50{,}6 - 18{,}6 = 32.
\]
\textbf{Đáp án của câu hỏi là B}
}
\end{ex}

\begin{ex} %[5D4V19]
Một hình tam giác có độ dài ba cạnh lần lượt là: $12{,}5\text{ cm}$; $15{,}8\text{ cm}$ và $17{,}5\text{ cm}$. Chu vi của hình tam giác đó là
\choice
{$44{,}8\text{ cm}$}
{\True $45{,}8\text{ cm}$}
{$46\text{ cm}$}
{$45{,}5\text{ cm}$}
\loigiai{
Chu vi hình tam giác là:
\[
(12{,}5 + 17{,}5) + 15{,}8 = 30 + 15{,}8 = 45{,}8\text{ (cm)}.
\]
\textbf{Đáp án của câu hỏi là B}
}
\end{ex}

\begin{ex} %[5D4V19]
Mẹ đi chợ mua $1{,}5\text{ kg}$ thịt bò, $2{,}25\text{ kg}$ cá và $1{,}75\text{ kg}$ rau củ. Tổng khối lượng mẹ đã mua là
\choice
{$5\text{ kg}$}
{\True $5{,}5\text{ kg}$}
{$5{,}25\text{ kg}$}
{$6\text{ kg}$}
\loigiai{
Tổng khối lượng là:
\[
1{,}5 + (2{,}25 + 1{,}75) = 1{,}5 + 4 = 5{,}5\text{ (kg)}.
\]
\textbf{Đáp án của câu hỏi là B}
}
\end{ex}

\begin{ex} %[5D4V19]
Tìm số tự nhiên $x$ thỏa mãn:
\[
12{,}4 + 15{,}8 < x < 14{,}5 + 15{,}2.
\]
\choice
{$x = 28$}
{\True $x = 29$}
{$x = 30$}
{$x \in \{28, 29\}$}
\loigiai{
Ta có:
\[
12{,}4 + 15{,}8 = 28{,}2; \quad 14{,}5 + 15{,}2 = 29{,}7.
\]
Bất đẳng thức trở thành:
\[
28{,}2 < x < 29{,}7.
\]
Số tự nhiên duy nhất thỏa mãn là $x = 29$.\\
\textbf{Đáp án của câu hỏi là B}
}
\end{ex}

\begin{ex} %[5D4V19]
Tính tổng của dãy số thập phân cách đều sau:
\[
S = 1{,}2 + 2{,}3 + 3{,}4 + 4{,}5 + 5{,}6 + 6{,}7 + 7{,}8 + 8{,}9.
\]
\choice
{$38{,}4$}
{\True $40{,}4$}
{$40$}
{$41{,}4$}
\loigiai{
Ghép cặp số đầu và số cuối:
\begin{align*}
S &= (1{,}2 + 8{,}9) + (2{,}3 + 7{,}8) + (3{,}4 + 6{,}7) + (4{,}5 + 5{,}6) \\
&= 10{,}1 + 10{,}1 + 10{,}1 + 10{,}1 = 10{,}1 \times 4 = 40{,}4.
\end{align*}
\textbf{Đáp án của câu hỏi là B}
}
\end{ex}

\begin{ex} %[5D4V19]
Một người thợ dệt, ngày thứ nhất dệt được $28{,}5\text{ m}$ vải. Ngày thứ hai dệt được nhiều hơn ngày thứ nhất $3{,}7\text{ m}$ vải. Ngày thứ ba dệt được nhiều hơn ngày thứ hai $1{,}8\text{ m}$ vải. Cả ba ngày người thợ đó dệt được số mét vải là
\choice
{$92{,}5\text{ m}$}
{\True $94{,}7\text{ m}$}
{$95\text{ m}$}
{$93{,}7\text{ m}$}
\loigiai{
Ngày thứ hai dệt được: $28{,}5 + 3{,}7 = 32{,}2\text{ (m)}$.\\
Ngày thứ ba dệt được: $32{,}2 + 1{,}8 = 34\text{ (m)}$.\\
Cả ba ngày dệt được: $28{,}5 + 32{,}2 + 34 = 94{,}7\text{ (m)}$.\\
\textbf{Đáp án của câu hỏi là B}
}
\end{ex}

\begin{ex} %[5D4V19]
Tổng của hai số thập phân là $48{,}6$. Nếu số thứ nhất tăng thêm $5{,}4$ và số thứ hai tăng thêm $2{,}8$ thì tổng mới là
\choice
{$55{,}8$}
{\True $56{,}8$}
{$56{,}2$}
{$57{,}8$}
\loigiai{
Tổng mới là:
\[
48{,}6 + 5{,}4 + 2{,}8 = 54 + 2{,}8 = 56{,}8.
\]
\textbf{Đáp án của câu hỏi là B}
}
\end{ex}

\begin{ex} %[5D4V19]
Số thích hợp điền vào dấu hỏi chấm để $15{,}72 + ? = 20$ là
\choice
{$4{,}38$}
{\True $4{,}28$}
{$5{,}28$}
{$4{,}18$}
\loigiai{
Số cần điền là:
\[
20 - 15{,}72 = 4{,}28.
\]
\textbf{Đáp án của câu hỏi là B}
}
\end{ex}

\begin{ex} %[5D4V19]
Ba túi gạo có khối lượng: túi 1 nặng $12{,}5\text{ kg}$; túi 2 nặng $14{,}75\text{ kg}$; túi 3 nặng $16{,}25\text{ kg}$. Tổng khối lượng của cả ba túi gạo là
\choice
{$42{,}5\text{ kg}$}
{\True $43{,}5\text{ kg}$}
{$44\text{ kg}$}
{$43{,}25\text{ kg}$}
\loigiai{
Tổng khối lượng là:
\[
12{,}5 + (14{,}75 + 16{,}25) = 12{,}5 + 31 = 43{,}5\text{ (kg)}.
\]
\textbf{Đáp án của câu hỏi là B}
}
\end{ex}

\begin{ex} %[5D4V19]
Tìm $x$ biết: $x - 18{,}45 = 12{,}55 + 4{,}8$.
\choice
{$34{,}8$}
{\True $35{,}8$}
{$36$}
{$35{,}2$}
\loigiai{
Tính vế phải: $12{,}55 + 4{,}8 = 17{,}35$.\\
Suy ra:
\[
x = 17{,}35 + 18{,}45 = 35{,}8.
\]
\textbf{Đáp án của câu hỏi là B}
}
\end{ex}

\begin{ex} %[5D4V19]
Một hình chữ nhật có chiều rộng là $8{,}75\text{ m}$, chiều dài hơn chiều rộng $3{,}5\text{ m}$. Chu vi hình chữ nhật đó là
\choice
{$21\text{ m}$}
{\True $42\text{ m}$}
{$38{,}5\text{ m}$}
{$40\text{ m}$}
\loigiai{
Chiều dài hình chữ nhật là:
\[
8{,}75 + 3{,}5 = 12{,}25\text{ (m)}.
\]
Nửa chu vi là:
\[
8{,}75 + 12{,}25 = 21\text{ (m)}.
\]
Chu vi hình chữ nhật là:
\[
21 \times 2 = 42\text{ (m)}.
\]
\textbf{Đáp án của câu hỏi là B}
}
\end{ex}

\textbf{Phần 2. Trắc nghiệm đúng sai}
\setcounterpageref{ex}{0}

\begin{ex} %[5D4V19]
Cho các phép tính cộng sau:
$A = 12{,}45 + 8{,}55$; $B = 7{,}8 + 13{,}2$; $C = 5{,}64 + 14{,}46$
Xét tính đúng sai của các khẳng định sau:
\choiceTF
{\True Giá trị của biểu thức $A$ là $21$.}
{\True Giá trị của biểu thức $B$ là $21$.}
{Giá trị của biểu thức $C$ lớn hơn giá trị của biểu thức $A$.}
{\True Tổng $A + B + C$ có giá trị bằng $62{,}1$.}
\loigiai{
\begin{itemize}
    \item $A = 12{,}45 + 8{,}55 = 21$. Khẳng định a \textbf{Đúng}.
    \item $B = 7{,}8 + 13{,}2 = 21$. Khẳng định b \textbf{Đúng}.
    \item $C = 5{,}64 + 14{,}46 = 20{,}1 < 21 = A$. Khẳng định c \textbf{Sai}.
    \item $A + B + C = 21 + 21 + 20{,}1 = 62{,}1$. Khẳng định d \textbf{Đúng}.
\end{itemize}
\begin{itemize}
    \item a) \textbf{Đúng}.
    \item b) \textbf{Đúng}.
    \item c) \textbf{Sai}.
    \item d) \textbf{Đúng}.
\end{itemize}
}
\end{ex}

\begin{ex} %[5D4V19]
Một đoàn xe tải chở hàng cứu trợ lũ lụt gồm ba xe:
- Xe 1 chở được $3{,}85\text{ tấn}$
- Xe 2 chở được nhiều hơn xe 1 là $0{,}65\text{ tấn}$
- Xe 3 chở được nhiều hơn xe 2 là $0{,}5\text{ tấn}$
Xét tính đúng sai của các khẳng định sau:
\choiceTF
{\True Khối lượng hàng xe 2 chở được là $4{,}5\text{ tấn}$.}
{\True Khối lượng hàng xe 3 chở được là $5\text{ tấn}$.}
{Cả ba xe chở được tất cả $14\text{ tấn}$ hàng.}
{\True Khối lượng hàng xe 3 chở nhiều hơn xe 1 là $1{,}15\text{ tấn}$.}
\loigiai{
\begin{itemize}
    \item Xe 2: $3{,}85 + 0{,}65 = 4{,}5\text{ tấn}$. Khẳng định a \textbf{Đúng}.
    \item Xe 3: $4{,}5 + 0{,}5 = 5\text{ tấn}$. Khẳng định b \textbf{Đúng}.
    \item Cả 3 xe: $3{,}85 + 4{,}5 + 5 = 13{,}35\text{ tấn} 
e 14\text{ tấn}$. Khẳng định c \textbf{Sai}.
    \item Xe 3 hơn xe 1: $5 - 3{,}85 = 1{,}15\text{ tấn}$. Khẳng định d \textbf{Đúng}.
\end{itemize}
\begin{itemize}
    \item a) \textbf{Đúng}.
    \item b) \textbf{Đúng}.
    \item c) \textbf{Sai}: Tổng là $13{,}35\text{ tấn}$.
    \item d) \textbf{Đúng}.
\end{itemize}
}
\end{ex}

\textbf{Phần 3. Trả lời ngắn}
\setcounterpageref{ex}{0}

\begin{ex} %[5D4V19]
Tính nhanh giá trị của biểu thức sau:
\[
P = 18{,}25 + 14{,}6 + 11{,}75 + 5{,}4.
\]
\shortans{50}
\loigiai{
\textbf{Đáp số:} $50$\\
Nhóm các số hạng:
\[
P = (18{,}25 + 11{,}75) + (14{,}6 + 5{,}4) = 30 + 20 = 50.
\]
}
\end{ex}

\begin{ex} %[5D4V19]
Một hình chữ nhật có chiều rộng là $14{,}5\text{ m}$, chiều dài là $20{,}5\text{ m}$. Hãy tính chu vi của hình chữ nhật đó theo đơn vị mét.
\shortans{70}
\loigiai{
\textbf{Đáp số:} $70$\\
Chu vi hình chữ nhật là:
\[
(14{,}5 + 20{,}5) \times 2 = 35 \times 2 = 70\text{ (m)}.
\]
}
\end{ex}

\begin{ex} %[5D4V19]
Khi cộng một số tự nhiên với $25{,}6$, bạn Lan đã quên viết dấu phẩy của số $25{,}6$ nên kết quả tìm được là $348$. Hãy tìm kết quả đúng của phép cộng đó. (Nhập dạng số thập phân, ví dụ $117.6$).
\shortans{117.6}
\loigiai{
\textbf{Đáp số:} $117.6$\\
Khi quên dấu phẩy, số $25{,}6$ biến thành $256$.\\
Số tự nhiên đó là: $348 - 256 = 92$.\\
Kết quả đúng của phép cộng là: $92 + 25{,}6 = 117{,}6$.
}
\end{ex}

\begin{ex} %[5D4V19]
Ba bạn An, Bình, Cường đi cân sức khỏe: An nặng $34{,}8\text{ kg}$; Bình nặng hơn An $2{,}4\text{ kg}$; Cường nặng hơn Bình $1{,}8\text{ kg}$. Hỏi cả ba bạn nặng tất cả bao nhiêu ki-lô-gam? (Nhập dưới dạng số thập phân, ví dụ $111$).
\shortans{111}
\loigiai{
\textbf{Đáp số:} $111$\\
Bình nặng: $34{,}8 + 2{,}4 = 37{,}2\text{ (kg)}$.\\
Cường nặng: $37{,}2 + 1{,}8 = 39\text{ (kg)}$.\\
Cả ba bạn nặng: $34{,}8 + 37{,}2 + 39 = 72 + 39 = 111\text{ (kg)}$.
}
\end{ex}

\begin{ex} %[5D4V19]
Tìm số tự nhiên $n$ biết:
\[
n - 15{,}4 = 24{,}6.
\]
\shortans{40}
\loigiai{
\textbf{Đáp số:} $40$\\
Ta có: $n = 24{,}6 + 15{,}4 = 40$.
}
\end{ex}

\begin{ex} %[5D4V19]
Tính tổng của bốn số thập phân sau: $12{,}5$; $27{,}8$; $17{,}5$ và $22{,}2$.
\shortans{80}
\loigiai{
\textbf{Đáp số:} $80$\\
Nhóm các số hạng:
\[
(12{,}5 + 17{,}5) + (27{,}8 + 22{,}2) = 30 + 50 = 80.
\]
}
\end{ex}
